\documentclass[../main]{subfiles}
\begin{document}
\ifSubfilesClassLoaded{\setcounter{page}{27}}{}
\section{Общая (общественная) собственность}
\label{sec:05_obshhaya_sobstvennost}

Препятствием для движения к коммунизму является частная собственность. Первая работа, где Маркс серьёзно и основательно задумывается о коммунизме – это «Экономическо-философские рукописи 1844 года». Глава о коммунизме находится во второй рукописи сразу после главы о частной собственности. Она начинается рассуждениями о частной собственности, о труде как о сущности частной собственности и о понимании частной собственности и труда у различных представителей социализма и коммунизма как идейного течения. Коммунисты выступают против частной собственности (в том числе и до Маркса) на средства производства, поскольку она делает людей рабами других людей и порождают все ужасы капиталистического общества.

Частная и личная собственность - это не одно и то же. Предметы непосредственного пользования и потребления, принадлежащие их владельцу, по форме находятся в частной собственности, точно так же, как и заводы, фабрики, магазины, корпорации, большие земельные участки. Это обстоятельство затуманивает суть дела. Потому, когда антикоммунисты слышат об уничтожении частной собственности, самые примитивные из них думают, что речь идёт об обобществлении абсолютно всех вещей и начинают бояться за свои личные вещи. В этом пункте мы можем их успокоить. Предметы личного потребления коммунистов интересуют вовсе не в том смысле, чтобы кого-то их лишить, а лишь в том, чтобы все были ими обеспечены в достаточной для полноценной и разносторонне-богатой жизни мере. Экономически эти предметы находятся не в частной, а в личной собственности.

Уничтожение частной собственности не отменяет личной собственности (частная собственность — это отношение между людьми по поводу вещей, а не вещи), поскольку потребление с необходимостью должно носить личный характер. Наоборот, уничтожение частной собственности обеспечивает доступ всех членов общества к предметам, необходимым для личного потребления, а там, где это более рационально, переводит потребление в общественные формы. Например, организация потребительских кооперативов позволяет повысить эффективность использования вещей и уменьшает потребность в их производстве и хранении, а значит также и затраты общественного времени на это. Так, организация простого потребительского кооператива по использованию предметов, которые используются не так часто, позволяет не захламлять предметами пространство и максимально эффективно использовать каждый из них вплоть до их морального и физического износа. При этом предмет переходит из рук в руки, и каждый участник кооператива может получить полезный эффект от него. Примеры потребительских кооперативов описывались в литературе много раз. Это легко проделывается с самыми разными предметами, которые перестают принадлежать кому-то лично и при этом используются всеми в должной, для удовлетворения потребностей мере. Коммунизм противостоит здесь личной собственности совсем не с той стороны, с которой опасаются его противники. Общественные формы потребления направлены не на то, чтобы у кого-то что-то отнять, а на то, чтобы дать доступ к предметам потребления всем.

Что касается таких общественных форм быта, как совместное использование стиральных машин, совместное приготовление и доставка пищи, совместное использования уборочной техники, то тут экономия ресурсов, и человеческого времени (и прямо, и косвенно – как экономия времени на производство ресурсов) при том же уровне потребления может быть очень большой. В каждый момент времени она пропорциональна возможностям техники. Но даже сейчас, например, приготовление каши для одного ребёнка в домашних условиях и приготовление каши на 200 детей в условиях детского сада или молочной кухни, оснащенных современной техникой, занимает одно и то же время.

То же самое касается и приготовления пищи для взрослых. Общественное потребление позволяет использовать более эффективные технологии, неприменимые в индивидуальных домашних условиях, с намного более высоким уровнем автоматизации всех процессов. При этом можно лучше вести контроль качества, что в итоге выгодно всем участникам процесса. Движение в этом направлении существуют и при капитализме. Сам капитализм на свой лад производит и стандартизирует систему общепита (Макдональдс, например, предлагает одинаковую картошку фри во всех частях света), прачечных и т.д. В условиях капитализма они всё равно не используются общественным способом. Формы совместного потребления постоянно наталкиваются на противодействие в товарном производстве, а производство, как известно, определяет потребление.

Все тяготы быта, в котором задействовано множество людей, можно заменить передовым производством, основанным на современных технологиях; мешающий друг другу и прохожим, и создающий пробки личный транспорт заменяется на хорошо развитый общественный, и так далее во всех сферах человеческой жизни, где это пошло бы на пользу всем. Например, в современных больших городах, где развит общественный транспорт, те, кто им пользуется, особенно в часы пик, добираются до нужного места более эффективно. В Москве дети, которые ездят в школу на метро в основном приезжают вовремя, те же, кого подвозят на личном автомобиле, как правило, опаздывают, поскольку попытки использования общественного ограниченного городского пространства частным способом, оказываются менее эффективны, чем общественное использование. Этим объясняются и призывы городских властей во время снегопадов отказаться от использования личного транспорта и пересесть на общественный, чтобы минимизировать аварийные ситуации на дорогах.

Товарное производство культивирует индивидуальные формы потребления. Даже при всей их нерациональности и неэффективности, они являются условием постоянного спроса на товары на рынке, условием движения капитала. Именно в этом смысле частная собственность соприкасается с личной и влияет на неё.

Частная собственность, или, точнее будет сказать, частная собственность на средства производства, это, главным образом, возможность извлечения прибавочного продукта (прибавочной стоимости, в форме прибыли или в форме ренты), то есть эксплуатация других людей.

Коммунисты говорят об обобществлении частной, а не личной собственности – средств общественного производства, которые сейчас используются не на благо всему обществу, а служат эксплуатации человека человеком. Частная собственность — это не вещи, а общественные отношения при посредстве вещей. Её сущность составляет эксплуатация.

Представители коммунизма до Маркса понимали уничтожение частной собственности как распространение её на всё общество, превращение всех членов общества в рабочих и уравнивание всех как в отношении труда, так и в отношении потребления. Частная собственность рассматривается многими представителями домарксового коммунизма только как вещественное богатство, которым можно владеть всем вместе. Например, Прудон мыслил частную собственность только как капитал, который подлежит уничтожению. Ближе к Марксу были те теоретики, которые считали источником частной собственности и её губительного влияния на человека определённый характер труда, «труд нивелированный, раздробленный и поэтому несвободный», однако они сводили такой труд к определённым видам труда. «Фурье, который, подобно физиократам, опять-таки считает земледельческий труд по меньшей мере наилучшим видом труда, тогда как по Сен-Симону, наоборот, суть дела заключается в промышленном труде как таковом, и он в соответствии с этим домогается безраздельного господства промышленников и улучшения положения рабочих. И, наконец, коммунизм есть положительное выражение упразднения частной собственности; на первых порах он выступает как всеобщая частная собственность». «Положительное», в данном случае означает «с сохранением всех богатств культуры и деятельных способностей человека, развитых в обществе товарного производства».

Понимание отрицания частной собственности как распространение её на всё общество подвергается Марксом беспощадной критике. Он показывает, что такой коммунизм не преодолевает частную собственность, а является лишь теоретическим завершением частной собственности как общественного отношения. Он отмечает, что господство собственности, как совокупности вещей, над таким коммунизмом, как идейно-теоретическим положением и общественным движением, ещё очень велико. Общественные отношения мыслятся ещё только под частнособственническими определениями. Такой коммунизм «стремится уничтожить все то, чем, на началах частной собственности, не могут обладать все; он хочет насильственно абстрагироваться от таланта и т. д. Непосредственное физическое обладание представляется ему единственной целью жизни и существования; категория рабочего не отменяется, а распространяется на всех людей; отношение частной собственности остается отношением всего общества к миру вещей; наконец, это движение, стремящееся противопоставить частной собственности всеобщую частную собственность, выражается в совершенно животной форме, когда оно противопоставляет браку (являющемуся, действительно, некоторой формой исключительной частной собственности) общность жен, где, следовательно, женщина становится общественной и всеобщей собственностью».

Для Маркса идея общности жен – это отношение к человеку как к вещи, которая показывает и отношение человека к миру вещей, созданных человеком для человека «Подобно тому как женщина переходит тут от брака ко всеобщей проституции, так и весь мир богатства, т. е. предметной сущности человека, переходит от исключительного брака с частным собственником к универсальной проституции со всем обществом». Маркс называет такой коммунизм грубым. Под богатством здесь понимается вся совокупность вещей, созданных людьми для людей, в которых воплощены как деятельные способности человеческого общества, так и потребности, которые оно выработало.

Теория грубого коммунизма, отрицающая личность человека, только абсолютизирует в мысленном отражении реальную тенденцию общества частной собственности. Там, где богатство человеческих отношений низводится до отношений частной собственности, личность низводится до экономического агента частной собственности. Собственно, личность этих агентов (как пролетария, так и капиталиста) для движения частной собственности абсолютно безразлична. Пролетарий интересует капиталиста не как личность, а исключительно как продавец рабочей силы определённого качества, капиталист пролетария – исключительно как её покупатель. То же происходит в актах купли- продажи любых других товаров. Все остальные личностные качества экономических агентов к движению частной собственности отношения не имеют. Если люди всё время своей жизни не выходят за рамки этих отношений, их личность начинает сводиться только к этим определениям.

Грубый коммунизм предполагает такое упразднение частной собственности, которое не было бы подлинным освоением ее. Это видно из абстрактного отрицания достижений культуры, рассмотрения их как излишних только потому, что они не нужны пролетарию как пролетарию и (по выражению Маркса) стремления к неестественной простоте бедного, грубого и не имеющего потребностей человека, который не может преодолеть частную собственность, так как ещё не дорос до неё.

Лозунг «Всё отнять и поделить», характерный для современных представителей псевдокоммунизма, стоит ещё ниже. Они даже не дошли до обобществления хотя бы в том виде, в котором её предлагали представители грубого коммунизма – обобществления нищеты.

Грубым коммунистическим теориям Маркс противопоставляет коммунизм как положительное упразднение частной собственности. Такой коммунизм, как исторический    процесс,    сохраняет    богатство  всего   предшествующего исторического развития в виде деятельных способностей, потребностей, материальной и духовной культуры человечества. Коммунизм делает всё это достоянием каждого члена общества. Как исторический процесс он носит сознательный, а не стихийный характер, поскольку обобществление предполагает сознательное отношение членов общества к производству и воспроизводству своей общественной жизни. Коммунизм предполагает освоение объективных законов общественного развития, которые человечество может поставить себе на службу точно так же, как законы физики, химии и биологии в современной промышленности.

Вся предшествующая история человечества порождает условия для такого коммунизма, который преодолевает частную собственность, но в этом преодолении сохраняет все богатство предшествующего развития. В процессе всей предшествующей истории формируется и соответствующее сознание – теория коммунизма. Коммунизм понимается как процесс уничтожения частной собственности: «Для уничтожения идеи частной собственности вполне достаточно идеи коммунизма. Для уничтожения же частной собственности в реальной действительности требуется действительное коммунистическое действие».

Следует отметить один важный, на наш взгляд, момент. Речь идёт об употреблении терминов «социализм» и «коммунизм», привычное значение которых устоялось не сразу. В цитируемых выше Рукописях 1844 года социализмом ещё называется общество, где частная собственность не снимается, а уже снята. В то время как коммунизм понимается как процесс (акт) отрицания (снятия) частной собственности. Именно поэтому здесь он не рассматривается как форма общества, а лишь как энергический принцип ближайшего будущего.

Это обстоятельство позволяет пролить свет на то, почему именно такие вопросы были подняты и именно так изложены Марксом и Энгельсом всего через 3 года после написания «Рукописей» в «Манифесте коммунистической партии». Чётче прослеживаются корни основного в Манифесте положения: «Но современная буржуазная частная собственность есть последнее и самое полное выражение такого производства и присвоения продуктов, которое держится на классовых антагонизмах, на эксплуатации одних другими. В этом смысле коммунисты могут выразить свою теорию одним положением: уничтожение частной собственности».

Это снятие, как уже отмечалось, предполагает сохранение всего культурного богатства, деятельных способностей и потребностей, развитых на основе частной собственности. Оно осуществляется не для самого себя, а для существования коммунистического общества с соответствующей собственностью и способом общения.

Каким бы ни было коммунистическое общество, человек как природное существо должен жить природой и присваивать себе природу: всё, чем пользуется и что употребляет человек, сделано из природного материала, при помощи природных процессов, поставленных на службу человеку, или является самими этими природными процессами. Через 20 лет после написания «Манифеста» во Введении из «Экономических рукописей 1857–1859» годов Маркс писал: «Всякое производство есть присвоение индивидуумом предметов природы в пределах определённой общественной формы и посредством неё. В этом смысле будет тавтологией сказать, что собственность (присвоение) есть условие производства. Смешно, однако, делать отсюда прыжок к определённой форме собственности, например, к частной собственности (что к тому же предполагает в качестве условия противоположную форму — отсутствие собственности)».

Коммунистическому обществу соответствует общая, то есть общественная форма собственности на все общественные средства производства и, соответственно, общее общественное управление всем общественным производством как единым целым. Общественному характеру производства, развитому на базе частной собственности, должен соответствовать общественный характер присвоения продуктов этого производства. Маркс и Энгельс писали об этом десятки раз, в разных своих работах. Пожалуй – это единственное положение, которое запомнили современные коммунисты, даже если они не понимают его значения.

Для начала этого исторического движения в наиболее развитых странах Маркс и Энгельс предложили следующие меры: «1. Экспроприация земельной собственности и обращение земельной ренты на покрытие государственных расходов.

2. Высокий прогрессивный налог. 3. Отмена права наследования. 4. Конфискация имущества всех эмигрантов и мятежников. 5. Централизация кредита в руках государства посредством национального банка с государственным капиталом и с исключительной монополией. 6. Централизация всего транспорта в руках государства. 7. Увеличение числа государственных фабрик, орудий производства, расчистка под пашню и улучшение земель по общему плану. 8. Одинаковая обязательность труда для всех, учреждение промышленных армий, в особенности для земледелия. 9. Соединение земледелия с промышленностью, содействие постепенному устранению различия между городом и деревней. 10. Общественное и бесплатное воспитание всех детей. Устранение фабричного труда детей в современной его форме».

Как мы видим, с точки зрения поставленной задачи по уничтожению частной собственности, эти меры выглядят не столь уж радикально, но общо, как и должно быть в манифестах. Манифесты должны быть ясны, понятны, и обозначать основное направление развития. Тем более, сами авторы утверждали, что меры зависят от конкретно-исторических условий, а потому могут быть иными. О мерах мы специально будем говорить в одном из следующих выпусков, через призму опыта социализма ХХ века.

Теперь же нужно сказать ещё пару слов о том, что такое частная собственность, что именно коммунисты ставят себе целью уничтожить с сохранением всего положительного содержания, чтобы «на место старого буржуазного общества с его классами и классовыми противоположностями» пришла «ассоциация, в которой свободное развитие каждого является условием свободного развития всех».

Частная собственность — это продукт несвободного труда и, в то же время, средство производства несвободного характера труда: такого труда, когда и процесс, и продукт труда не принадлежит непосредственно коллективу трудящихся людей, труда, который сам по себе обезразличен для трудящихся, а является лишь средством по достижению внешних целей (денег). Сущность частной собственности – в таком характере труда. Он доходит до пределов своего развития в капиталистическом обществе развитого товарного производства, благодаря доведению до крайней степени общественного разделения труда, превращающего продукты труда в товары (в товарном производстве).

Общественное разделение труда долгое время способствовало развитию человечества. Оно повышало производительность труда и создавало материальную базу для того, чтобы передать машиноподобные функции, которые вынужден выполнять человек из-за неразвитости орудий, самим этим орудиям- машинам. Но оно не только объединяет людей, дескать, Вася - слесарь, а Петя - программист, каждый делает своё дело, а над ними ещё стоит специалист- менеджер, который управляет ими для общего блага. Такая цель не преследуется никем из участников процесса. Все люди при таком положении дел, оказываются связанными между собой внешним образом – через продукты своего труда. Они вынуждены обменивать свой продукт на тот, которого у них нет, оставаясь в процессе труда безразличными друг к другу. Не важно, кто и как делает этот продукт, важно, что этот продукт чужого труда нужен мне.

Общественное разделение труда прошло несколько крупных стадий, сыгравших положительную роль в развитии человечества. О них имеет смысл говорить в отдельной работе, специально посвященной этому вопросу. Здесь только важно отметить, что общественное развитие – продукт истории. Марксисты отметают с порога все предрассудки о биологически-природной приспособленности членов общества к тем или иным видам труда. Эти предрассудки не выдерживают даже поверхностной научной критики. Они не соответствуют не только тому, что известно об истории этого вопроса, но и той истории, которая совершается на наших глазах. Ежегодно появляется и исчезает множество новых профессий. Родители часто не понимают, в чём суть работы их детей и уж точно не могли передать к ней склонность по наследству вместе с генами.

С тех пор, как труд становится частичным, а продукты труда людей становятся товарами, то есть начинают производиться специально для обмена, эти продукты начинают господствовать над людьми. Вещи – продукты их собственного труда – и деньги, как представители всех вещей, приобретают над живыми людьми хорошо известную власть. А с тех пор, как товаром становится рабочая сила, такое господство, соответственно, приобретают люди, у которых эти деньги есть в достаточном количестве, чтобы купить средства производства и рабочую силу. Эти люди владеют живым человеческим трудом. Экономически они являются только представителями капитала. Их функция – присоединять к капиталу чужой труд, чтобы увеличилась стоимость капитала. В конце концов, это приводит к тому, что почти все члены общества так или иначе работают ради увеличения капитала – самовозрастания стоимости. «Только продукты самостоятельных, друг от друга независимых частных работ противостоят один другому как товары». Поэтому уничтожение частной собственности – это уничтожении товарного производства вообще, а не только капиталистического товарного производства.

Стоимость стремится к самовозрастанию. Это закон, управляющий обществом, где каждый прикован к своей частичной деятельности или, другими словами, общества развитого товарного производства. Законы современных государств – естественное юридическое выражение этого объективного закона производства, основанного на несвободном и раздробленном труде.

До определённого момента стремление капитала к самовозрастанию является движущей силой развития общества. Но оно вызывает к жизни общественные силы, которым становится тесно в рамках капитала. Использование их исключительно в целях увеличения капитала ведет к катастрофическим последствиям.

Управлять этими силами нужно, да и возможно на благо общества только всем вместе для развития всех, по общему разумному плану. А для этого необходимо освоить универсальные человеческие способности (овладеть универсальной чувственностью, мышлением и нравственностью), выйдя за рамки своей профессии, не позволяющие разобраться даже в «соседней» профессии, а уж тем более - в общественных процессах в целом. Поэтому уничтожение частной собственности равно уничтожению несвободного характера труда. Оно осуществляется по мере развития свободной коммунистической деятельности и превращение её в достояние всех людей.

Частной собственности коммунизм противопоставляет общую собственность, как форму совместного практического отношения людей к средствам производства и природным ресурсам. Общая (общественная) собственность обеспечивает наиболее эффективное использование производственных возможностей не только для общества в целом, но и для каждого его члена. Она предполагает активное участие и активную волю в отношении общественных ресурсов от каждого члена общества. Соответственно, если частная собственность – продукт разделения труда (раздробление его на всё более частные виды труда)\footnote{Именно потому что в СССР общественное разделение труда не было преодолено (тенденция к его преодолению не стала преобладающей), а пролетариат воспроизводился как пролетариат, то есть как класса буржуазного общества, произошла полная реставрация господства частной собственности.} – и является этим разделением, то общая собственность предполагает, как уже отмечалось выше, обобществление труда, устранение его разделения. Общественная собственность предполагает свободную человеческую деятельность по меркам истины, добра и красоты. Последнее и составляет реальную экономическую основу для реального, а не только формального обобществления и реального совпадения частных и общественных интересов всех членов общества.

Под общей или общественной собственностью не следует понимать собственность коллективных хозяйств и кооперативов, или иных групп людей, которые производят товары и торгуют ими со всем остальным миром. Они не выходят за пределы частной собственности, так как не выходят за пределы товарного производства. С обществом члены таких хозяйств связаны как товаровладельцы и товаропроизводители. Сам характер товарного производства определяет и отношения в этих коллективах, так как ставит производство стоимости во главу угла, делает прибыль целью работы коллектива.

Зародыш общей собственности также не следует искать, например, в собственности семьи. В ХІХ веке некоторые социалисты представляли себе дело таким образом, будто общественная собственность может получиться путём «разрастания» семейной собственности до масштабов всего общества. Общество представлялось в виде большой семьи с характерными для семьи способами распределения.    Но семья – воплощение частного характера присвоения продуктов труда. Она сама в её современной форме есть продукт развития частной собственности. В той мере, в которой в семье осуществляется производство и воспроизводство человека как рабочей силы (человек как экономического агента общества товарного производства), она целиком вписана и полностью воспроизводится по законам движения частной собственности. В той мере, в которой она воспроизводит всех своих членов и является основной потребительской «ячейкой», она целиком и полностью определяется капиталистическим производством. Это экономически предполагается как в форме заработной платы, так и в форме прибыли. Более того, современная семья как хозяйственная единицы не только не исключает, но и прямо предполагает эксплуатацию человека человеком, что несовместимо с коммунизмом. Потому собственность семьи нельзя считать основой общественной собственности – собственности, характерной для коммунизма.

И уж тем более, коммунистическую собственность, которая могла бы быть узловой точкой движения к коммунизму, нельзя увидеть во всякого рода изолированных сообществах, не связанных с общественным производством (например, религиозные коммуны, основанные на натуральном хозяйстве).

Но где же тогда можно увидеть общую собственность характерную для коммунизма? Очевидно, что только там, где есть свободная человеческая деятельность в своей непосредственно-коллективной форме. Такая деятельность, взятая со стороны использования и распоряжения ресурсами, со стороны присвоения при помощи общественной формы и является наличным бытием общей собственности.

В предыдущей главе мы привели лишь некоторые примеры такой коммунистической деятельности в социалистическом и капиталистическом обществе. Если брать эти и другие примеры со стороны волеизъявления по отношению к предмету труда, средствам труда и продукту, речь будет идти о формах общественной собственности и общественном характере присвоения.

Например, общедоступная база данных по своей сути уже не может находиться не в общественной собственности, не говоря уже об общей системе управления экономикой. То же относится и к свободному ПО, и к разнообразным сервисам, предполагающим совместное производство и использование. Существуют также серверы, по совместному использованию ресурсов, которые неэффективно и невыгодно использовать на началах частной собственности, например, общие библиотеки и хранилища для фильмов и музыки, предполагающие, что участники и получают, и раздают продукты на этих сервисах (но пример с ПО и базами лучше, так как здесь речь идёт скорее о производстве, чем о потреблении).

Из сказанного выше видно, что общественная собственность рождается из реального движения производства даже там, где ещё нет формального обобществления средств производства, в условиях капитализма. Формальное обобществление – взятие политической власти пролетариатом и овладением им средствами производства – даст огромный простор для развития реальной общественной собственности. Заметим, что там, где есть реальная, а не только формальная общественная собственность, исчезает почва для того, чтобы пытаться присвоить себе часть общественных ресурсов, или наоборот разбазаривать их, небрежно к ним относиться как к «ничьим».

При формальном обобществлении требуется внешнее регулирование хотя бы в форме моральной сознательности (на самом деле в виде закона и системы принуждения, обеспечивающей его исполнение) в отношении к средствам, предмету и продуктам труда, исключающее подобные вещи. При реальном обобществлении потребность в таком регулировании отпадает постольку, поскольку культура в отношении к средствам производства становится простым условием их использования.
\end{document}
